\section{Cross-Cutting Process Environment and Additives}

The firing environment is part of the material history, not background metadata. In a stored abstract, otherwise related Al-containing pellets fired in alumina and platinum crucibles are compared using composition, relative density, ionic conductivity, and air-exposure readouts \cite{xia2016ionic}. The record supports the statement that crucible choice and exposure were experimental variables in that study; it does not license a universal crucible prescription. A separate title-level record explicitly concerns alumina-crucible contamination associated with extra lithium carbonate \cite{lan2019investigation}, so it is retained as a topic anchor without importing a causal mechanism.

Lithium inventory must be split temporally. Starting excess is set in the batch, whereas lithium activity during firing can be influenced by atmosphere, cover material, geometry, and dwell. A title-only paper in the collection is framed around manipulation of a Li2O atmosphere for dense LLZO \cite{huang2019manipulating}. An abstract-bearing cover-powder study reports LiOH-containing cover material together with density, total conductivity, activation energy, and grain-boundary composition \cite{zhang2019densification}. Another abstract screens several lithium-containing, lower-cost cover-powder candidates for compensation of firing-stage lithium loss \cite{huang2019searching}. These records map related controls, but only a matched design could separate starting inventory from boundary supply.

Exposure continues after firing. A stored abstract presents lithium--proton exchange as a variable to be controlled for reproducible LLZO processing and performance \cite{rosen2021controlling}. That association makes humidity, exposure time, storage state, surface treatment, and delay before impedance measurement candidates for the minimum record. Their absence from a checked source is recorded as NR; it is not interpreted as zero exposure.

Additives create a second attribution problem. The title of a W-containing garnet paper explicitly links a sintering aid with microstructure and conductivity topics \cite{luo2019influence}, but its local title-only provenance prevents any detailed effect claim. Abstract-bearing studies associate SiO2, CuO, AlN, or Li5AlO4 additions with combinations of densification, boundary or phase interpretations, transport, and cell readouts \cite{zhang2020enhanced,zhang2019effects,zhang2022high,zhou2024li5alo4}. Because additive identity, lithium supply, reaction products, density, and boundary transport may move together, no single abstract establishes which mediator is necessary or sufficient across compositions.

The practical comparison unit is consequently not ``additive present.'' It is the tuple \{parent powder, nominal and measured additive content, mixing, green-body state, firing schedule and atmosphere, cover\slash crucible, phase assemblage, density basis, boundary characterization, conductivity type and temperature\}. The uniform-densification abstract shows why geometry and spatial sampling also belong to that tuple \cite{guo2023uniform}. With the tuple complete, a study may support a within-paper causal test; with only an additive name and an endpoint, it remains an association.
